\section{从字符到带类型程序：预处理与编译前端}
\label{sec:frontend}

源文件看似已经“写明”程序，但字符本身并不知道 \texttt{clipped\_dot} 是函数、\texttt{a[i]}
是左值、\texttt{-8} 的负号是运算符，或 \texttt{putint} 应接收一个整数。本节沿着字符、预处理
记号、语法结构和带类型 AST 四层表示，说明这些事实如何逐层成立。所谓\emph{编译前端}
（compiler frontend）在这里指面向源语言、产出带语义结构的部分；它不是“所有生成 LLVM IR
以前发生的命令”的同义词。

\subsection{先澄清：SysY 语义载体与 C 观察载体}

规范 SysY 文件不写 \texttt{\#include}，因为 SysY 运行时函数由编译器按语言约定识别；课程的
SysY 语法也没有定义 C 预处理器。为了满足对完整语言处理系统的观察要求，又不污染规范源，实验
另有一个 C17 观察载体。它的首行和所包含头文件中的关键指令如下：

\begin{verbatim}
#include "case_config.h"

/* case_config.h */
#ifndef PREFLIGHT_CASE_CONFIG_H
#define PREFLIGHT_CASE_CONFIG_H
#define VECTOR_LENGTH 8
#define TERM_MIN (-8)
#define TERM_MAX 12
int getint(void);
void putint(int value);
void putch(int value);
#endif
\end{verbatim}

\texttt{case\_config.h} 为 Clang 提供 \texttt{getint/putint/putch} 的 C 函数声明；三个对象式宏
提供数组长度和两个截断界；包含卫士（include guard）用 \texttt{\#ifndef}/\texttt{\#define}
保证同一翻译单元内重复包含时声明只进入记号流一次。这里必须诚实说明：当前 C 载体不是直接
包含 \texttt{.sy}，而是手工镜像函数体。它与规范 SysY 的精确语法差异还包括输入数组形参的
\texttt{const} 限定、\texttt{main(void)}、用 \texttt{else if} 简写嵌套条件、引入
\texttt{result} 临时变量以及 Doxygen 注释。这些差异不改变式~\eqref{eq:clipped-dot} 的可观察
函数，但手工镜像会带来漂移风险；共享已知答案、结构检查和代码审阅是当前控制措施，不能把它
冒充“单一源文件生成”。

这种安排把三类陈述分开：（1）SysY 文档规定的语言/运行时事实；（2）ISO C 与 Clang 文档规定
的预处理和 C 前端事实；（3）这个兼容子集在指定 Clang 版本中的观察。库存 Clang 接受观察用
翻译单元，并不能证明它接受所有合法 SysY 或拒绝所有非法 SysY。

\subsection{翻译阶段：规范顺序与实现融合}

C 标准用若干\emph{翻译阶段}（translation phases）描述“仿佛按此顺序”发生的处理：源文件字符
映射、反斜杠换行拼接、识别预处理记号并把注释替换为空格、执行预处理指令并展开宏、处理字符
与字符串字面量、相邻字符串拼接，最后把预处理记号转换为翻译记号并完成语法/语义分析；各翻译
单元随后交给链接阶段\cite{front-c-n1570}。这一语义顺序不要求实现为多趟文件扫描。GNU CPP
明确说明最初几步会为了性能一并完成\cite{front-gcc-cpp-initial}；Clang 中词法器向预处理器持续
提供记号，预处理器再把展开后的流交给解析器，故“先完整词法分析、再完整预处理”同样不是准确
的进程时间线\cite{front-clang-internals}。

这里需要区分\emph{预处理记号}（preprocessing token）与稍后带语言意义的记号。预处理阶段先
识别标识符、预处理数、字符串/字符字面量和标点，使宏替换按记号而非任意字符串进行；一些
预处理数即使不是合法的 C 数值常量，也可先作为预处理记号存在。宏展开完成后，解析器看到的是
结果流，而不是“宏调用 AST”。这解释了为什么源中 \texttt{VECTOR\_LENGTH} 若是宏，会在 AST
中表现为整数常量（通常保留拼写/展开位置以支持诊断），却没有一个运行时读取“宏变量”的动作。

\subsection{包含、宏与条件编译分别改变什么}

\paragraph{包含。}
\texttt{\#include} 常被说成“粘贴头文件文本”，但更精确的描述是：实现定位所选文件，把它作为
输入继续扫描、执行其中指令和宏，然后回到原文件。查找结果受引号/尖括号形式、\texttt{-I}、
系统目录、资源目录、目标 sysroot 等影响，因此可复现记录应保存实际搜索路径，而不是只保存
\texttt{\#include "case\_config.h"}。\texttt{clang -E -H} 或详细模式能显示选择了哪个物理文件；
预处理输出中的行标记则在展开文本与原文件位置间建立诊断映射\cite{front-gcc-cpp-include}。

\paragraph{宏。}
对象式宏把一个标识记号替换为记号序列；函数式宏还要收集实参、按规则预展开、替换并重新扫描。
字符串化 \texttt{\#} 和记号粘合 \texttt{\#\#} 对实参展开有特殊规则，正在展开的宏还需暂时
禁用，以免无效递归。因而宏不是无条件文本替换，更不是类型化函数调用\cite{front-gcc-cpp-macro}。
在观察载体里，\texttt{VECTOR\_LENGTH} 被替换为 \texttt{8}；
\texttt{\#define TERM\_MIN (-8)} 的括号也属于替换序列，而 \texttt{-8} 仍在后续
词法/语法层成为一元负号作用于常量 8。预处理器不知道该表达式最终类型是不是 \texttt{int}。

\paragraph{条件编译。}
\texttt{\#if}/\texttt{\#ifdef} 决定哪些预处理记号进入后续流。被 \texttt{\#if 0} 排除的分支
无需构成合法 C 语法，因为解析器根本不会接收它；不过条件指令自身仍须被预处理器正确配对。
这与运行时 \texttt{if} 根本不同：案例中 \texttt{if (term < lo)} 会进入 AST、IR 和目标控制流，
而未选中的预处理分支在解析前已经消失。条件编译的能力适合平台适配，也可能令构建配置成为
程序语义的隐含输入。本例实际观察的是头文件包含卫士：第一次包含后宏
\texttt{PREFLIGHT\_CASE\_CONFIG\_H} 已定义，若再次包含，受保护正文便被排除。

可用以下停止点建立可检查证据：
\begin{verbatim}
clang -std=c17 -Ipreflight/include -E \
      preflight/src/clipped_dot.c -o clipped_dot.i
clang -std=c17 -Ipreflight/include -E -dM \
      preflight/src/clipped_dot.c > macros.txt
clang -fsyntax-only -Xclang -dump-tokens \
      -std=c17 -Ipreflight/include \
      preflight/src/clipped_dot.c \
      2> tokens.txt
\end{verbatim}
\texttt{.i} 是预处理记号流的一种可读投影，并非预处理器内部状态的完整序列化；项目的自动取证
执行第一条命令并确认三个宏名已消失、常量值已出现。后两条是补充观察方法：\texttt{-dM} 还会
包含预定义和头文件宏，不能误称为“案例自己定义的宏列表”。

\subsection{词法：从拼写到记号}

\emph{词法分析}（lexical analysis）把字符组织成带种类、拼写和源位置的记号。对循环片段
\texttt{term = a[i] * b[i];}，解析器接收的核心次序可概括为：
\[
\begin{aligned}
 &\textit{identifier(term)},=,\textit{identifier(a)},[,\textit{identifier(i)},],*,\\
 &\textit{identifier(b)},[,\textit{identifier(i)},],;.
\end{aligned}
\]
空格与注释主要用于分隔和诊断定位，不成为普通 AST 节点；\texttt{while} 被识别为关键字而非
可由符号表重新解释的普通标识符；\texttt{<}、\texttt{>}、\texttt{=} 的区别在这一层确定。

负常量尤其适合展示层次边界。在规范 SysY 调用的实参 \texttt{-8}（C 载体中则来自宏
\texttt{TERM\_MIN}）里，通常没有单个“负八字面量”记号，而是 \texttt{-} 与整数常量
\texttt{8} 两个记号。语法分析把前者解释为一元运算，
语义分析再确定其操作数和结果的整数类型，常量求值可以得到 $-8$。到了 LLVM IR，
\texttt{i32 -8} 又可以作为一个常量操作数出现。不同层的文本相似或不同，都不能单凭拼写判断
它们是同一种对象。

词法器还须处理“最长匹配”等规则：\texttt{else if} 是两个关键字，\texttt{<=} 是一个标点记号，
而 \texttt{< =} 是两个；\texttt{//} 开始的注释在翻译早期被空白替代。错误的字符序列会产生
词法诊断，但“\texttt{term} 是否已声明”不是词法问题——那需要名字查找。

\subsection{语法：记号如何形成结构}

\emph{语法分析}（parsing）依据文法把线性记号流组织成声明、语句和表达式结构。Clang 使用
递归下降解析器，并在识别产生式时调用 Sema；它并非先构造一棵完整的具体语法树，再启动一个
完全分离的语义遍历\cite{front-clang-internals}。因此“词法→语法→语义”仍是责任上的良好分层，
但不应被叙述为三遍彼此隔绝的实现。

以 \texttt{acc = acc + term;} 为例，运算符优先级使右侧先形成加法表达式，整体再成为赋值表达式；
左结合规则使连续加法按规定分组。\texttt{a[i] * b[i]} 先形成两个下标表达式，再形成乘法节点。
\texttt{else} 与最近尚未配对的 \texttt{if} 结合，使乘积恰好落入“小于下界”“大于上界”“区间
内”三种路径之一。AST 可简化表示为：

\begin{verbatim}
WhileStmt
|- condition: BinaryOperator '<'
|  |- DeclRefExpr i : int
|  `- DeclRefExpr n : int
`- CompoundStmt
   |- VarDecl term : int
   |  `- BinaryOperator '*'
   |     |- ArraySubscriptExpr a[i] : int lvalue
   |     `- ArraySubscriptExpr b[i] : int lvalue
   |- IfStmt
   |  |- BinaryOperator '<' (term, lo)
   |  `- ... else IfStmt ('>' (term, hi))
   |- BinaryOperator '=' (acc, '+' (acc, term))
   `- BinaryOperator '=' (i, '+' (i, 1))
\end{verbatim}

这棵摘录是为阅读而压缩的结构图，不冒充某次 \texttt{-ast-dump} 的逐字输出。真实 Clang AST
通常还含源范围、隐式转换和类型信息；开发者文本/JSON dump 是观察接口，不是跨版本稳定的交换
格式\cite{front-clang-ast}。

\subsection{Sema：名字、作用域、类型与隐式操作}

\emph{语义分析}（semantic analysis）回答“结构是否在该语言中有意义”。Clang 将这组职责集中
在 Sema 库：声明与名字查找、作用域、类型检查、隐式转换、常量表达式约束以及带源位置的诊断。
Sema 在解析期间构造/完善 AST，所以最终 AST 不只是语法括号树，而是带类型的源语言语义图景。

\paragraph{符号与作用域。}
\emph{符号表}（symbol table）把标识符出现与声明实体联系；\emph{作用域}（scope）规定某个声明
在哪些源位置可见。顶层的 \texttt{clipped\_dot} 和 \texttt{main} 是函数实体；函数体内形参
\texttt{a,b,n,lo,hi} 与局部变量 \texttt{acc,i,term} 分属相应块作用域。\texttt{term} 只在循环
体的复合语句内可见，而 \texttt{i} 的生存期覆盖整个函数体。每一个 \texttt{DeclRefExpr i}
都应解析到同一个局部声明，而不是仅凭字符串相等猜测。若在声明 \texttt{i} 之前使用它，字符和
语法都可能完全合法，名字查找仍应拒绝程序。

\paragraph{类型与值类别。}
\texttt{a} 的源形参写作 \texttt{int a[]}，在函数参数位置调整为指向整数的参数；在
\texttt{a[i]} 中，整数索引选择一个 \texttt{int} 元素。该表达式先是可定位对象的\emph{左值}
（lvalue），乘法需要它的值，于是 AST 中可见隐式的左值到右值转换。把数组实参传给
\texttt{clipped\_dot} 时则发生数组到指针转换（array-to-pointer conversion），不是复制八个
整数。\texttt{n}、\texttt{lo}、\texttt{hi} 按值传入。Clang AST 有意接近源语言，通常以
\texttt{ImplicitCastExpr} 暴露这些虽未写出却由语言要求的操作\cite{front-clang-ast}。

表~\ref{tab:semantic-nodes} 将几个关键节点的前端义务汇总起来。

\begin{table}[t]
\caption{案例节点的语义分析结果与后续义务。}
\label{tab:semantic-nodes}
\centering
\small
\begin{tabular}{p{3.0cm}p{4.4cm}p{4.2cm}}
\hline
节点 & 前端确立的事实 & 交给后续层的义务 \\
\hline
\texttt{a[i]} & \texttt{a} 指向 \texttt{int}，\texttt{i} 为整数，结果为 \texttt{int} 左值 & 形成正确元素地址并加载 32 位值 \\
\texttt{a[i]*b[i]} & 两操作数经值转换后为 \texttt{int}，结果为 \texttt{int} & 按源语言整数语义乘法；案例证明无溢出 \\
\texttt{term < lo} & 两侧为 \texttt{int}，结果可作条件 & 建立真假控制流且只执行所选路径 \\
\texttt{clipped\_dot(...)} & 解析到五参数函数；两地址、三整数实参与形参匹配 & 形成 ABI 正确的调用或语义等价内联 \\
\texttt{putch(10)} & 运行时声明解析为接收一个整数、无返回值 & 把 10 传给外部调用并保留副作用顺序 \\
\hline
\end{tabular}
\end{table}

\paragraph{常量规则。}
\emph{常量表达式}（constant expression）是在规定上下文可于编译期求值、并满足语言额外限制的
表达式。SysY 要求数组各维长度为可编译期求值的非负整数；规范源中的字面量 \texttt{8}
显然满足约束。C 载体里的 \texttt{VECTOR\_LENGTH} 则先由预处理器替换为同一字面量。
同理，规范源的 \texttt{-8} 在词法上是 \texttt{-} 和 \texttt{8}，经语义检查与常量求值得到
整数 $-8$；C 载体的 \texttt{TERM\_MIN} 先展开为带括号的相同记号序列。常量求值不等于宏
展开：前者在
类型和语言语义已经介入后求值，后者在解析前替换记号。

数组初始化也需要语义处理。八个初值分别转换并写入对应元素；若列表短于数组长度，SysY 规则会
对余下元素作零初始化，而本例刻意给满，避免把隐式补零混入点积结果。局部变量
\texttt{acc}、\texttt{i}、\texttt{term} 都在读取前有明确定义，从而不以未初始化值作为跨层契约。

\subsection{诊断是失败的语义产物}

错误消息不是编译器外的附属打印，而是前端在无法建立表示契约时产生的用户界面。一个受控变异
可以把用户调用改为
\begin{verbatim}
clipped_dot(a, b, n, TERM_MIN);  /* missing TERM_MAX */
\end{verbatim}
词法器仍能产生全部记号，解析器也能形成四实参调用的语法；Sema 在把调用解析到五形参函数后，
才能指出实参数量不足，并把诊断定位到调用及函数声明。另一个变异 \texttt{a[1.5]} 在 C 中可被
解析为下标结构，但下标类型约束失败。这样的变异应位于隔离的诊断实验中，不应污染规范案例。

诊断同时说明为什么 AST 是“已解释的源程序”而非简单解析树：对于合法调用，AST 的被调实体、
函数类型和隐式转换均已确定；对于非法调用，编译器可能为错误恢复构造部分节点，但后续代码生成
没有义务把它当作合法程序。一个稳健的报告应保存命令、工具版本、退出状态和关键诊断，而不对
措辞逐字符断言，因为诊断文本不是稳定 ABI。

\subsection{AST 的终点也是下一层的起点}

到带类型 AST 为止，前端已经知道：两个数组元素是 32 位整数；比较按有符号整数意义进行；
\texttt{while} 有一个回到条件的结构；运行时调用具有特定函数类型。但它尚未决定循环值放在哪个
物理寄存器、数组在栈帧何处、调用地址是多少，也不必永久保存宏展开的高层意图。Clang IRGen
消费 AST，把结构化控制转成基本块和分支，把左值访问转成地址与加载/存储，把函数调用转成带
类型 IR 调用。LLVM 的经典设计正是让语言相关语义先由前端消化，再进入可供跨语言优化的带类型
SSA 表示\cite{front-lattner-adve}。

这是一处有损但受控的抽象下降：源位置和调试元数据可以保留部分出处，然而后续优化器不应被假定
能恢复 \texttt{else if} 的原始拼写、宏定义边界或变量为何命名为 \texttt{term}。下一章因此不再
问“这行源代码被翻成哪一行 IR”，而问更稳健的问题：控制依赖、数据依赖、内存对象和调用副作用
分别由哪些 IR 对象承载。表示改变了，契约没有中断。
